% Où limiter : client, bord du réseau, passerelle, service ; la clé à chaque étage.
\documentclass[border=6pt]{standalone}
\usepackage{sd}
\begin{document}
\begin{tikzpicture}[x=1mm, y=1mm, note/.append style={font=\sffamily\small},
    etage/.style={soft=#1, rounded corners=3pt, minimum width=34mm, minimum height=10mm, font=\sffamily\small\bfseries}]
  \foreach \x/\c/\t in {0/Dslate/client, 44/Dteal/{bord (CDN, WAF)}, 88/Dblue/passerelle d'API, 132/Dpurple/service} {
    \node[etage=\c] (n\x) at (\x,40) {\t};
  }
  \draw[arr=Dink] (n0) -- (n44);
  \draw[arr=Dink] (n44) -- (n88);
  \draw[arr=Dink] (n88) -- (n132);
  \foreach \x/\cle/\quoi in {
      0/{respecte Retry-After}/{recul exponentiel, gigue,\\limitation adaptative},
      44/{clé : adresse IP, préfixe}/{volume par IP, inondations,\\robots ; grossier},
      88/{clé : clé d'API, utilisateur}/{quota du forfait, rafales,\\429 et en-têtes RateLimit},
      132/{clé : locataire, ressource}/{requêtes en cours, opérations\\coûteuses, délestage}} {
    \node[note, anchor=north, font=\sffamily\small\bfseries] at (\x,32) {\cle};
    \node[note, anchor=north, align=center] at (\x,26) {\quoi};
  }
  \draw[arr=Dred] (88,47) to[bend right=20] node[lblc=Dred, above] {429 + Retry-After} (0,47);
\end{tikzpicture}
\end{document}
